\section*{ID9354: Definition of ΔT𝜙}
\paragraph{Metadata.}
PiID: 2798394004152 | RTEID: RTE:P36323.6593:T3.3016 | UUID: 9b0e4122-d89e-53c9-9dfa-2dd8f6d80662

\paragraph{Canonical Statement.}
n audit: compiled corpus certificate. Rocq audit: compiled corpus certificate. Dependencies and Test Handle Dependencies [ID 0017] Node registry\_id\_0018 Support Titles Regularized Index Raising; Second Trawinistic Derivative; Normalized Laplacian Test Handle Observed diffusion or dispersion behavior that cannot be matched by any regularized second-order operator built from the Trawinistic connection would challenge the definition. Canonical Equation Links \$∇\$ 𝑎 T =𝑔𝑎𝑏reg \$∇T\$ 𝑏 \$∇T 𝑎∇\$ 𝑎 T 𝜙 \$ΔT𝜙=∇T 𝑎∇\$ 𝑎 T 𝜙 [ID 0019] Trawinistic d’Alembertian TZPID ID0019 UUID 6d74c597-ddfc-52d2-adb1-f8d00b173427 Type Operator Scale wave-operator LOC Classification Working routing: QC174.5 wave propagation; QA431 hyperbolic operators arXiv Classification Primary: gr-qc Cross-list: hep-th, math-ph Canonical Statement The Trawinistic d’Alembertian is the wave operator induced by the regularized inverse metric an

\paragraph{Canonical Equation.}
\[
ΔT𝜙=∇T 𝑎∇
\]

\paragraph{Dictionary.}
Definition of ΔT𝜙 — ΔT𝜙=∇T 𝑎∇

\paragraph{Encyclopedia.}
Definition of ΔT𝜙 is an algebraic definition, written ΔT𝜙=∇T 𝑎∇. It is an algebraic definition expressing the quantity directly in terms of the variables on its right-hand side. It is built from T, defining ΔT𝜙. Within the Trawin Zero Point registry it serves as one element of the framework's mathematical narrative.
